% THE Q-MODEL — Reference Monograph
% Version 1.0 (revised) — White-paper reference edition
% Build: pdflatex TQM_v1_0_Monograph.tex (twice, then makeindex if an index is added)

\documentclass[11pt,twoside]{book}
\usepackage[utf8]{inputenc}
\usepackage[T1]{fontenc}
\usepackage{amsmath,amssymb,amsthm}
\usepackage[margin=2.7cm]{geometry}
\usepackage{booktabs}
\usepackage{array}
\usepackage{longtable}
\usepackage{lmodern}
\usepackage[hidelinks]{hyperref}
\usepackage{xcolor}

\definecolor{caveat}{HTML}{7A2E2E}

% Theorem-like environments
\newtheorem{theorem}{Theorem}[chapter]
\newtheorem{proposition}[theorem]{Proposition}
\newtheorem{lemma}[theorem]{Lemma}
\newtheorem{corollary}[theorem]{Corollary}
\newtheorem{definition}[theorem]{Definition}
\newtheorem{axiom}[theorem]{Axiom}
\newtheorem{remark}[theorem]{Remark}
\newtheorem{nogo}[theorem]{No-Go Theorem}

\newcommand{\derived}{\textbf{DERIVED}}
\newcommand{\realunder}{\textbf{REAL-UNDERIVED}}
\newcommand{\drawn}{\textbf{DRAWN}}

\title{\bfseries THE Q-MODEL — From Q to Cosmology\\[3mm]
\large A Reference Monograph on Structure, Complexity and Random Actualization}
\author{Fabrice Wieser}
\date{Version 1.0 (revised) \quad·\quad 15 August 2026}

\begin{document}

\frontmatter
\maketitle

\begin{center}
\fcolorbox{caveat}{white}{%
\begin{minipage}{0.92\textwidth}
\color{caveat}
\textbf{Publication status: READY\_FOR\_WHITEPAPER — NOT\_READY\_FOR\_JOURNAL.}\\[2mm]
This reference monograph accompanies the white paper \emph{THE Q-MODEL — From Q to
Cosmology} and is released for archival on Zenodo. It is \emph{not} submitted for
peer-reviewed journal publication. The central derivation claim (Einstein recovery from
Q-events) remains \emph{logical, not mathematical}: the metric and the BDG action are
imported (proven but not TQM-derived), $G=\ell^2c^3/\hbar$ is dimensional analysis, and no
unique sharp prediction yet discriminates TQM from SM$+\Lambda$CDM. Every claim marked
\textbf{verified} is backed by an executable xUnit test in
\texttt{TQM.Tests/ResearchXC} (or \texttt{ResearchQG}); see Part~V and Appendix~A.
\end{minipage}}
\end{center}

\vspace{3mm}

\chapter*{Abstract}
\addcontentsline{toc}{chapter}{Abstract}
\noindent
This monograph is the reference edition of THE Q-MODEL (TQM), a theory of structure and
content that compresses observable physics to four primitives — the individuation
principle $Q$, Random Actualization, the scale triad $(\ell,\tau,\hbar)$, and a single
continuous nonlinearity parameter $M^2$ — and holds that the \emph{form} of every structure
is derivable from these primitives while the \emph{content} is realized by contingent draw.
The monograph presents, in self-contained form: the formal primitives and dynamical system
(the graph Laplacian $L_Q$ and the Schr\"odinger equation); the verified continuum-limit
chain from $L_Q$ to the flat Laplacian, the d'Alembertian, the Laplace--Beltrami operator,
the curved-space Schr\"odinger equation, and the Einstein tensor; the derivation hierarchy
(gauge structure, spatial dimensionality, the multiplicity bound, the abundance law,
phase-gradient gravity); the three-category taxonomy; the conditional no-go theorems; the
quantitative predictions; and the complete inventory of executable verification results.

\tableofcontents

\mainmatter

%=============================================================================
\part{Foundations}
%=============================================================================

\chapter{The Primitives}
\label{ch:primitives}

\section{Statement of the theory}

THE Q-MODEL rests on a single organizational claim, the \textbf{structure/content split}:

\begin{center}
\emph{Structure is derivable; content is realized.}
\end{center}

The \emph{form} of observable structure (topology, symmetry, distribution shape, lower
bounds) is provable from a minimal primitive set. The \emph{content} (specific masses,
couplings, multiplicities, angles) is a draw from a contingent ensemble and is therefore
classified rather than derived.

\section{The four primitives}

\begin{definition}[Primitive set]
The theory admits exactly four primitives, and no others:
\begin{enumerate}
\item \textbf{$Q$} — the principle of individuation (ontology): the irreducible act by
which a discrete event comes to be individuated from nothing.
\item \textbf{Random Actualization} — genuine ontological chance: given $Q$, the realized
event locations are random within the causal structure. (Assumption \textbf{A-03}.)
\item \textbf{$(\ell,\tau,\hbar)$} — the irreducible physical triple: a spacetime scale
$\ell$, a clock $\tau$, and an action $\hbar$ (unit conventions).
\item \textbf{$M^2$} — the single continuous nonlinearity parameter (the nonlinearity
regime of the dynamics).
\end{enumerate}
\end{definition}

No gauge-group, multiplicity, or hidden-dimension primitive is permitted. The primitives
have two layers that meet at $Q$: an \emph{ontology} layer (underwriting the derivation of
structure) and a \emph{dynamical} layer (underwriting the dynamical system of
Chapter~\ref{ch:dynamics}).

\begin{remark}
$Q$ individuates; Random Actualization realizes; $(\ell,\tau,\hbar)$ fixes units; $M^2$
sets the single continuous parameter that the derivation hierarchy pins to $M^2\approx5$
(Chapter~\ref{ch:complexity}).
\end{remark}

\section{The quantum postulates (dynamical layer)}

\begin{axiom}[Q exists]
Topological charge quanta have position $x_i$ and phase $\theta_i$, and interact pairwise
with coupling $J_{ij}=f(|x_i-x_j|)$. (Assumed, irreducible.)
\end{axiom}

\begin{axiom}[Reversible dynamics]
The state norm $\|\psi\|^2$ is conserved; equivalently, evolution is unitary.
\end{axiom}

\begin{axiom}[Born rule]
$P=|\langle\phi|\psi\rangle|^2$, uniquely selected by additivity (Gleason's theorem, 1957).
\end{axiom}

\begin{axiom}[Measurement]
A measurement yields one definite outcome (the collapse axiom).
\end{axiom}

Postulates 1--2 derive the Hilbert space and the Schr\"odinger equation
(Chapter~\ref{ch:dynamics}); postulates 3--4 close the interpretation.

%=============================================================================

\chapter{Formalization of $Q$ and Random Actualization}
\label{ch:formalization}

This chapter records the formal status of the primitives, as established by the
formalization audits (executable; see Part~V).

\section{Status of $Q$}

\begin{center}
\begin{tabular}{lll}
\toprule
Component & Status & Content \\
\midrule
Object & \textbf{Formalized} & a discrete event (individuated entity) \\
State space & \textbf{Formalized} & positions $x_i$ and phases $\theta_i$ \\
Operations & \textbf{Formalized} & individuation, coupling $J_{ij}$ \\
Configuration space & \textbf{Partial} & graph of events, no continuum limit native \\
Dynamics & \textbf{Partial} & $L_Q$ generator (postulate 1) \\
Symmetries & \textbf{Partial} & $S^1$ phase, $Z_2$, permutation $S_n$ \\
Continuum limit & \textbf{Partial} & verified for $L_Q$ (Part~II) \\
Measure & \textbf{Missing} & no measure space from primitives alone \\
\bottomrule
\end{tabular}
\end{center}

\textbf{Classification:} $Q$ is \emph{partially formalized} — object, state space, and
operations exist; the measure and the action functional remain as postulates.

\section{Status of Random Actualization}

Random Actualization is \textbf{A-03}, an assumption, not a derived object. The
formalization audit found:

\begin{center}
\begin{tabular}{lll}
\toprule
Component & Status & Content \\
\midrule
Random variable & \textbf{Formalized} & the draw of event locations \\
Generator & \textbf{Formalized} & ontological chance \\
$(\Omega,\mathcal F,P)$ & \textbf{Partial} & ensemble unspecified \\
Ensemble measure & \textbf{Partial} & no explicit measure \\
\bottomrule
\end{tabular}
\end{center}

\textbf{Classification:} the random variable and generator are formalized; the probability
space $(\Omega,\mathcal F,P)$ and the ensemble measure are only partially formalized.

\begin{remark}
The hostile reviewer's charge that ``Random Actualization is verbalism'' is therefore
\emph{partially valid}: it is an \emph{admitted} primitive (A-03), not a derived object,
and the paper states this explicitly.
\end{remark}

%=============================================================================

\chapter{The Dynamical System}
\label{ch:dynamics}

\section{The graph Laplacian $L_Q$}

\begin{definition}[Graph Laplacian]
For the $Q$-interaction graph with adjacency $A$ and degree matrix $D$,
\begin{equation}
L_Q = D - A .
\end{equation}
$L_Q$ is real, symmetric, positive semi-definite, and has zero row sums. Its eigenvectors
form an orthonormal basis — the \textbf{Hilbert space} of the theory.
\end{definition}

\section{The tight-binding identity (TQM-142)}

\begin{proposition}[Tight-binding identity]
On a 1D chain,
\begin{equation}
(L_Q)_{ij} = 2\delta_{ij} - \delta_{i,j+1} - \delta_{i,j-1},
\end{equation}
which is \emph{identical} to the tight-binding Hamiltonian with on-site energy
$\varepsilon=2t$: $H = t L_Q$. This is a mathematical \textbf{identity}, not an analogy.
\end{proposition}

\section{Schr\"odinger from reversibility (TQM-149--151)}

\begin{theorem}[Schr\"odinger equation from reversibility]
Consider linear evolution $\partial_t\psi = M\psi$ with conserved norm
$\partial_t\|\psi\|^2 = \psi^\dagger(M+M^\dagger)\psi = 0$. Then
\begin{equation}
M^\dagger = -M .
\end{equation}
For real $M$, this is $M^T=-M$. The simplest $2\times2$ antisymmetric matrix $J$ satisfies
$J^2=-I$ (so $J$ acts as the imaginary unit), and
\begin{equation}
M = J\otimes L_Q \quad\Longrightarrow\quad i\partial_t \psi = L_Q\,\psi,
\end{equation}
with unitary evolution
\begin{equation}
\psi(t) = e^{-iL_Qt}\,\psi(0).
\end{equation}
\end{theorem}

\section{Observables from the spectrum (TQM-145)}

\begin{center}
\begin{tabular}{ll}
\toprule
Observable & Definition \\
\midrule
effective mass & $m_{\rm eff} = 1/\lambda_1$ \\
energy & $E = \mathrm{tr}(L)$ \\
gap & $\Delta = \lambda_2-\lambda_1$ \\
correlation length & $\xi = 1/\sqrt{\lambda_1}$ \\
diffusion constant & $D = \lambda_1$ \\
channel capacity & $C = \log_2 N$ \\
\bottomrule
\end{tabular}
\end{center}

\begin{remark}
These are \emph{dimensional assignments}, not dynamical predictions; the hostile reviewer
correctly notes this is a scope boundary.
\end{remark}

\section{The ontology $\to$ physics bridge}

The bridge has two legs, both now tested:
\begin{enumerate}
\item $L_Q$ eigenvector basis $=$ Hilbert space (TQM-149); and
\item the causal-set continuum limit yields the metric (QG-001, XC006--012).
\end{enumerate}

%=============================================================================
\part{The Continuum-Limit Program (Verified)}
%=============================================================================

\chapter{$L_Q \to -\nabla^2$: The Flat Laplacian}
\label{ch:flat-laplacian}

\begin{theorem}[Discrete spectrum of the chain Laplacian]
The graph Laplacian of a 1D chain of $N$ sites has the exact spectrum
\begin{equation}
\lambda_k = \frac{1}{\Delta x^2}\left[2 - 2\cos\frac{\pi k}{N+1}\right],
\qquad k=1,\dots,N,
\end{equation}
with $\Delta x$ the lattice spacing. In the continuum limit $N\to\infty$, $\Delta x\to0$
with the length fixed,
\begin{equation}
\lambda_k \to (\pi k)^2 .
\end{equation}
\end{theorem}

\begin{proof}
The chain Laplacian is the tridiagonal Toeplitz matrix with diagonal $2/\Delta x^2$ and
off-diagonal $-1/\Delta x^2$; its eigenvectors are the discrete sine modes
$\sin(\pi k j/(N+1))$.
\end{proof}

\textbf{Verification.} \texttt{GraphLaplacianContinuumTests} builds chains with
$N=32,64,128,256$, diagonalizes $L_Q$ via MathNet.Numerics, and compares against the closed
form. The relative error \emph{decreases} with $N$ (second-order convergence, error
$\sim4\times$ per doubling), establishing that $L_Q$ converges to the flat Laplacian with
controlled error. \textbf{Result: PASSED.}

\chapter{BDG $\to \square$: The d'Alembertian}
\label{ch:dalembertian}

The causal-set d'Alembertian is the Benincasa--Dowker--Glaser (BDG) operator: a
directed, alternating sum over causal intervals of the sprinkling.

\begin{theorem}[BDG convergence]
The BDG operator converges to the continuum d'Alembertian $\square$ as the sprinkling
density grows, with leading truncation error $O(h^2)$.
\end{theorem}

\textbf{Verification.} \texttt{BDGOperatorContinuumTests} verifies convergence with a
tolerance $10^{-2}$ (the leading $O(h^2)$ error at $h=1/16$ is $\approx4.37\times10^{-3}$);
the error falls $\sim4\times$ per $h$-halving. \textbf{Result: PASSED.}

%=============================================================================

\chapter{The Quantum-Gravity Bridge}
\label{ch:bridge}

The two continuum chains are \emph{partially connected} — and their signatures are
incompatible.

\begin{proposition}[Two distinct operators]
\begin{enumerate}
\item $L_Q$ is an \emph{undirected, positive} graph Laplacian whose continuum limit is
$-\nabla^2$ (Riemannian, elliptic, positive semi-definite).
\item The BDG operator is a \emph{directed, alternating} causal-set operator whose
continuum limit is $\square$ (Lorentzian, hyperbolic, indefinite).
\end{enumerate}
\end{proposition}

\textbf{Verification.} \texttt{QuantumGravityBridgeTests} (three tests) establish:
\begin{enumerate}
\item $L_Q$ is positive semi-definite (minimum eigenvalue $\ge0$);
\item the BDG operator is indefinite (possesses negative eigenvalues);
\item the two are incompatible in signature — the graph Laplacian cannot serve as a
Lorentzian d'Alembertian (BdgUniquenessAnalyzer O3/O6 explicitly rejects it).
\end{enumerate}
\textbf{Result: PASSED (3/3).}

\begin{remark}
This is the single most important structural fact of the continuum program: TQM possesses
\emph{two} chain types, one Riemannian ($L_Q\to-\nabla^2$) and one Lorentzian
($\mathrm{BDG}\to\square$), which are \emph{disjoint} in signature. The bridge between
them is not yet closed natively.
\end{remark}

%=============================================================================

\chapter{The Weighted Laplacian and Laplace--Beltrami}
\label{ch:laplace-beltrami}

\section{The weighted Laplacian}

\begin{definition}[Weighted Laplacian]
Given the spatial coupling matrix $K_{ij}$, the weighted graph Laplacian is
\begin{equation}
L_W = D_K - K,
\qquad (D_K)_{ii} = \sum_j K_{ij}.
\end{equation}
This is the standard \emph{unnormalized} graph Laplacian (Belkin--Niyogi), valid for
uniform density.
\end{definition}

\begin{proposition}[Properties]
$L_W$ is symmetric, has zero row sum, is positive semi-definite, and reduces to the
unweighted Laplacian when $K=A$.
\end{proposition}

\textbf{Verification.} \texttt{WeightedLaplacianTests} (four tests) confirm symmetry, zero
row sum, positive semi-definiteness, and reduction to the unweighted case.
\textbf{Result: PASSED (4/4).}

\section{The Laplace--Beltrami approximation}

\begin{theorem}[Laplace--Beltrami from $L_W$]
$L_W$ approximates the Laplace--Beltrami operator $\Delta_g$ in the flat limit and
reproduces standard manifold examples.
\end{theorem}

\textbf{Verification.} \texttt{LaplaceBeltramiTests} (three tests):
\begin{enumerate}
\item \emph{Flat limit}: $L_W$ preserves the flat spectrum;
\item \emph{Variable weights}: non-uniform $K$ changes the spectrum;
\item \emph{Known manifold}: the $S^1$ cycle graph reproduces the circle spectrum.
\end{enumerate}
A subtle degeneracy was found and fixed by audit: the nonzero $S^1$ modes are two-fold
degenerate ($e^{\pm ik\theta}$), so in ascending order the first occurrence of mode $k$ is
at index $2k-1$, not $k$. \textbf{Result: PASSED (3/3).}

%=============================================================================

\chapter{The Curved-Space Schr\"odinger Equation}
\label{ch:curved-schrodinger}

\begin{theorem}[Curved Schr\"odinger equation]
The operator $L_W$ defines a curved-space Schr\"odinger equation
\begin{equation}
i\,\partial_t \psi = L_W \psi .
\end{equation}
Because $L_W$ is self-adjoint, the evolution is unitary and the norm is conserved:
\begin{equation}
\partial_t \|\psi\|^2 = 0 .
\end{equation}
\end{theorem}

\textbf{Verification.} \texttt{CurvedSchrodingerTests} (three tests): unitarity (norm
conservation), reduction to the flat Schr\"odinger equation, and the existence of a
curved operator. \textbf{Result: PASSED (3/3).}

\begin{remark}
This establishes the \emph{Schr\"odinger} side of the curved-space bridge: the weighted
Laplacian $L_W$ produces a well-defined, unitary curved Schr\"odinger equation. The
\emph{metric} that would make $L_W$ a genuine $\Delta_g$ remains external (Chapter~10).
\end{remark}

%=============================================================================

\chapter{The Einstein Tensor Chain}
\label{ch:einstein-tensor}

\section{Standard differential geometry}

\begin{definition}[Christoffel symbols]
\begin{equation}
\Gamma^\lambda_{\mu\nu} = \tfrac12\, g^{\lambda\sigma}
\left( \partial_\mu g_{\sigma\nu} + \partial_\nu g_{\sigma\mu} - \partial_\sigma g_{\mu\nu} \right).
\end{equation}
\end{definition}

\begin{definition}[Riemann curvature]
\begin{equation}
R^\rho_{\ \sigma\mu\nu} = \partial_\mu \Gamma^\rho_{\nu\sigma}
- \partial_\nu \Gamma^\rho_{\mu\sigma}
+ \Gamma^\rho_{\mu\lambda}\Gamma^\lambda_{\nu\sigma}
- \Gamma^\rho_{\nu\lambda}\Gamma^\lambda_{\mu\sigma}.
\end{equation}
\end{definition}

\begin{definition}[Ricci and Einstein tensors]
\begin{equation}
R_{\sigma\nu} = R^\rho_{\ \sigma\rho\nu}, \qquad
G_{\mu\nu} = R_{\mu\nu} - \tfrac12 R\, g_{\mu\nu}.
\end{equation}
\end{definition}

\section{Verified examples}

\textbf{Unit 2-sphere} $g=\mathrm{diag}(1,\sin^2\theta)$ at $\theta=\pi/4$:
\begin{equation}
\Gamma^\theta_{\phi\phi} = -\sin\theta\cos\theta = -0.5,
\qquad K = 1, \qquad R = 2K = 2, \qquad R_{\mu\nu}=K g_{\mu\nu}.
\end{equation}

\textbf{Unit 3-sphere} $g=\mathrm{diag}(1,\sin^2\chi,\sin^2\chi\sin^2\theta)$ at
$\chi=\theta=\pi/4$: $R_{\mu\nu}=2g_{\mu\nu}$, $R=6$, hence
\begin{equation}
G_{\mu\nu} = R_{\mu\nu} - \tfrac12 R g_{\mu\nu} = -g_{\mu\nu}
= -\mathrm{diag}(1,\tfrac12,\tfrac14).
\end{equation}

\begin{theorem}[2D Einstein tensor vanishes]
In two dimensions $G_{\mu\nu}\equiv0$ identically ($R_{\mu\nu}=\tfrac12 R g_{\mu\nu}$); a
non-trivial $G_{\mu\nu}$ requires dimension $\ge3$.
\end{theorem}

\textbf{Verification.} \texttt{EinsteinTensorTests} (four tests: metric$\to$Christoffel,
Christoffel$\to$Riemann, Riemann$\to$Ricci, Ricci$\to$Einstein) and
\texttt{EinsteinTensorIntegrationTests} (four tests on the integrated builder
\texttt{EinsteinTensorBuilder}). \textbf{Result: PASSED (8/8).}

\begin{remark}[Where the chain breaks in TQM]
The standard chain is fully functional. But TQM's own analyzers
(\texttt{QuantumGravityEmergenceAnalyzer}, \texttt{GrBridgeAnalyzer}) \emph{describe} the
chain as strings (\texttt{GeoStep}) and \emph{never compute} it. The chain breaks at
Step~1 — \emph{metric} $\to$ \emph{Christoffels} — because TQM possesses no native metric
field $g_{\mu\nu}(x)$: it is imported (Chapter~10).
\end{remark}

%=============================================================================

\chapter{Metric Emergence and Origin Closure}
\label{ch:metric-emergence}

\section{The causal distance structure}

\begin{definition}[Causal volume and distance]
For a causal interval of depth $D$ in $d$ dimensions, the causal volume is
\begin{equation}
N(D) \propto D^d,
\end{equation}
and the causal distance is
\begin{equation}
L = N^{1/d}.
\end{equation}
\end{definition}

\textbf{Verification.} \texttt{MetricGenerationTests} and \texttt{MetricEmergenceTests}
confirm: $L=N^{1/d}$ recovers the depth exactly ($d=2,3,4$; $D=1,\dots,8$); the dimension
is recovered from volume growth; and the resulting distance matrix satisfies all four
metric axioms (non-negativity, identity of indiscernibles, symmetry, triangle inequality).
\textbf{Result: PASSED.}

\section{The conformal factor from the counting measure}

\begin{theorem}[Conformal factor]
The counting measure fixes the volume element $\sqrt{|g|}=\rho$ (density $\rho$). For a
conformally flat ansatz $g=f\,\eta$,
\begin{equation}
\sqrt{|g|} = f^{\,d/2} = \rho \quad\Longrightarrow\quad f = \rho^{\,2/d}.
\end{equation}
\end{theorem}

\textbf{Verification.} \texttt{MetricEmergenceTests} confirms the round-trip
$f^{d/2}=\rho$ for $\rho\in\{0.5,1,2,8\}$, $d\in\{2,3,4\}$. The conformally flat candidate
$g=f\eta$ is constructible: constant $f$ gives $R=0$; $f=1+\tfrac12 x^2$ gives
$R=-0.6145$, matching the Liouville result $K=-\tfrac1{2f}\nabla^2\ln f$.
\textbf{Result: PASSED.}

\section{The conformal class: Malament's theorem}

\begin{theorem}[Malament 1977; Hawking--King--McCarthy 1976]
A bijection preserving the causal order $\prec$ between two past-and-future-distinguishing
spacetimes is a conformal isometry: the causal order determines the metric \emph{up to a
conformal factor}.
\end{theorem}

\textbf{Verification.} \texttt{ConformalStructureTests} and \texttt{MetricOriginTests}
establish: the causal order is a valid partial order whose boundary is the light cone; the
null structure is invariant under $g\to f g$ ($f>0$) but \emph{changed} by a non-conformal
rescale $g=\mathrm{diag}(-1,2)$; hence the light cone picks out \emph{exactly} the
conformal class. \textbf{Result: PASSED.}

\begin{theorem}[Metric origin closure]
The conformal-class ``gap'' is an \emph{imported theorem} (Malament), not a theory gap.
The metric origin is closed:
\begin{equation}
Q\text{-events} \to \text{causal order (native)}
\to \text{conformal class (proven)}
\to \text{conformal factor (native)}
\to g_{\mu\nu}.
\end{equation}
\end{theorem}

\textbf{Classification of the metric origin:}

\begin{center}
\begin{tabular}{lll}
\toprule
Link & Status & Nature \\
\midrule
Q-events & \textbf{NATIVE} & TQM-derived primitive \\
Causal order & \textbf{NATIVE} & precedence from individuation \\
Conformal class & \textbf{IMPORTED} & Malament/HKM — \emph{proven} \\
Conformal factor & \textbf{NATIVE} & counting measure, $f=\rho^{2/d}$ \\
Full $g_{\mu\nu}$ & \textbf{determined} & class $\times$ factor \\
\bottomrule
\end{tabular}
\end{center}

\begin{remark}
The metric origin is \emph{already solved}: importing a proven theorem is standard, not a
defect. The genuinely open item is the \emph{native} metric $\to$ operator coupling
(the ``G4'' gap): TQM does not produce $g_{\mu\nu}$ from Q-events, it imports it.
\end{remark}

%=============================================================================
\part{The Derivation Hierarchy}
%=============================================================================

\chapter{Complexity and Spatial Dimensionality}
\label{ch:complexity}

\section{The complexity functional}

The complexity argument is a \textbf{weighted six-component decomposition} of the
conditions for observers (\texttt{ComplexityOptimumAnalyzer.cs}), not a single scalar
variational functional.

\begin{center}
\begin{tabular}{lcc}
\toprule
Component & Weight & $M^2$ window \\
\midrule
Structure formation & 3.0 & 2--8 \\
Particle stability & 2.5 & 3--7 \\
Chemistry potential & 4.0 & 3--5 \\
Information capacity & 2.0 & 3--10 \\
Evolution potential & 1.5 & 2--6 \\
Observer viability & 5.0 & 4--6 \\
\bottomrule
\end{tabular}
\end{center}

\section{Spatial dimensionality $d=3+1$}

\begin{theorem}[Unique dimensionality]
$d=3+1$ is the \emph{unique} value satisfying simultaneously: Bertrand's theorem (closed
orbits only for $1/r^2$), Gauss's law ($F\propto1/r^{d-1}$, Coulomb only in 3D), knot
stability (codimension-2), Huygens (sharp propagation in odd $d$), and 2 GR polarizations.
$d=2+1$ fails gravity/atoms/knots; $d\ge4+1$ fails orbits and knots.
\end{theorem}

\begin{remark}
$M^2\approx5$ is the observer-viability \emph{intersection} of the component windows, not a
hand-inserted number. Confidence T-02 = 0.85; this is a window-intersection argument, not a
variational theorem. The analyzer records that the generation count $G\approx3$ is a
\emph{plateau}, not a peak.
\end{remark}

%=============================================================================

\chapter{Gauge Structure}
\label{ch:gauge}

\section{$U(1)$ — derived}

\begin{theorem}[$U(1)$ from the circle]
Phase lives on $S^1$. Its isometry group is
\begin{equation}
\mathrm{Aut}(S^1) = U(1), \qquad \pi_1(S^1) = \mathbb{Z},
\end{equation}
and the winding $\oint\nabla\theta\cdot dl = 2\pi n$ yields integer topological charge.
(Confidence 0.95.)
\end{theorem}

\section{$SU(2)$ and $SU(3)$ — real-underived structure}

\begin{center}
\begin{tabular}{lll}
\toprule
Factor & Status & Confidence \\
\midrule
$U(1)$ & \derived & 0.95 \\
$SU(2)$ & \realunder\ (emergent) & 0.70 \\
$SU(3)$ structure & \realunder\ (emergent) & 0.10 \\
color count $n=3$ & \drawn & — \\
\bottomrule
\end{tabular}
\end{center}

Binary winding $\{n=\pm1\}$ gives the minimal doublet; the Bloch sphere gives
$SO(3)=SU(2)/Z_2$; the spinor double cover gives $SU(2)$. The automorphism of the
$n$-defect moduli space satisfies $\mathrm{Aut}(C^n/S_n)\supseteq SU(n)$, deriving
$SU(3)$ \emph{structure} but not the count.

\begin{remark}
TQM derives the gauge-group \textbf{structure} (which group), \textbf{not} the gauge
\textbf{dynamics} (Maxwell / Yang--Mills actions are borrowed; the 8-gluon algebra is
assumption A-07). The count $n=3$ is unfixed (gauge-count no-go T-09, provisional 0.10).
\end{remark}

%=============================================================================

\chapter{Flavor and the Koide Relation}
\label{ch:flavor}

\section{The Yukawa operator}

The Yukawa operator is the overlap operator
\begin{equation}
Y_{ij} = \langle \mathrm{arch}_i \,|\,\mathrm{amplitude}\,|\,\mathrm{arch}_j\rangle .
\end{equation}
Its \emph{form} is derived; its \emph{spectrum} (the hierarchy $1{:}207{:}3478$) is
underived — set by the contingent architecture shapes.

\section{The Koide relation}

\begin{theorem}[Koide relation]
The charged-lepton pole masses satisfy
\begin{equation}
Q = \frac{\sum m}{(\sum\sqrt{m})^2} = 0.6666605,
\qquad \theta = 44.9997^\circ,
\end{equation}
equivalently $Q=2/3$ at $\theta=\arccos(1/\sqrt2)=45^\circ$. The relation is real
($\sim10^{-5}$ precision; Bayes factor vs.\ coincidence $\approx3.2\times10^4$), predictive
(1981 prediction $m_\tau=1776.97$ MeV, confirmed 1992), and RG-stable.
\end{theorem}

\begin{nogo}[Koide origin — T-08]
No symmetry ($S_3$), attractor, topology ($S^1$), or information-geometry selects $Q=2/3$.
(Confidence 0.70.)
\end{nogo}

The closure audit (Phase 159) exhausted seven routes:

\begin{center}
\begin{tabular}{ll}
\toprule
Route & Status \\
\midrule
Symmetry ($S_3$) & No-Go \\
Topology ($S^1$) & No-Go \\
Attractors & No-Go \\
Information Geometry & No-Go \\
Group Theory & No-Go \\
$m=3$ Closure & No-Go \\
Moduli Arguments & Blocked \\
\bottomrule
\end{tabular}
\end{center}

The Koide origin is a \emph{closed question} — real, predictive, RG-stable,
\emph{underivable} — the canonical \realunder\ structure.

\begin{nogo}[Neutrino-Koide — T-11, falsification]
Requiring $Q=2/3$ for neutrinos gives $Q_{\max}=0.585$ (normal) / $0.500$ (inverted), both
$<2/3$; the measured $\Delta m^2$ exclude all-lepton Koide. (Confidence 0.90.)
\end{nogo}

%=============================================================================

\chapter{Gravity and Emergent GR}
\label{ch:gravity}

\section{Newton's constant}

\begin{proposition}[Dimensional result]
$G = \ell^2 c^3/\hbar$ is a \emph{unit-consistency} result, not a derivation of gravity's
dynamics.
\end{proposition}

\section{The phase-gravity chain}

oscillation $\to$ phase $\to$ causal density $\to$ metric $\to$ curvature $\to$ Einstein.

\begin{theorem}[Leading-order recovery]
\begin{equation}
G_{\mu\nu} = 8\pi G_{\rm eff}\, T_{\mu\nu} + O(\ell_P^2 R^2),
\end{equation}
with Planck-scale corrections $\sim10^{-40}$. Newtonian $1/r^2$, lensing, redshift,
perihelion precession, and two gravitational-wave polarizations at speed $c$ are reproduced.
\end{theorem}

Two modifications beyond GR: (1) time-varying dark energy
\begin{equation}
\Lambda(t) = \frac{\alpha}{\sqrt{V(t)}},
\end{equation}
where $V(t)$ is the 4-volume of the past light cone; and (2) singularity resolution at
$r\sim\ell_P$ (maximum curvature $1/\ell_P^2$).

\begin{remark}[Honesty note]
The chain is \textbf{logical, not a new dynamical derivation}: ``same equations, same
predictions; no falsifiable difference'' in the weak field. The value is \emph{ontological}
(identifying the gravitational potential with the phase field); the physical content beyond
GR is confined to $\Lambda(t)$ and singularity resolution.
\end{remark}

\section{The radial acceleration relation}

\begin{theorem}[Zero-parameter RAR]
\begin{equation}
g_\dagger = \frac{c H_0}{2\pi} \approx 1.05\times10^{-10}\ \mathrm{m/s^2},
\end{equation}
matching the measured $a_0$. (The factor $2\pi$ is a numerical accident, Phase 144.)
\end{theorem}

%=============================================================================

\chapter{Cosmology}
\label{ch:cosmology}

\begin{itemize}
\item \textbf{Expansion} is an interpretation of FLRW (QG-081); every reinterpretation
collapses to FLRW (QG-080--089).
\item \textbf{Causal-set $\Lambda$}: $\Lambda\sim1/\sqrt{N}$ is a postdiction (Phase 140).
\item \textbf{Pantheon+SH0ES}: TQM and $\Lambda$CDM are indistinguishable at the current
signal $w(z)=-1+0.015(1+z)^{3/2}$ (DATA-001); a detectability audit (DATA-002) quantified
the power required for separation.
\item \textbf{Clusters}: Coma dynamical mass and ACCEPT X-ray gas fraction are reproduced.
\item \textbf{CMB}: an accepted partial computational layer; full $C_\ell$ solver deferred
(computational, not a theory gap).
\end{itemize}

%=============================================================================
\part{Classification, No-Gos, and Predictions}
%=============================================================================

\chapter{The Three-Category Taxonomy}
\label{ch:taxonomy}

\begin{definition}[Taxonomy]
\begin{center}
\begin{tabular}{ll}
\toprule
Category & Meaning \\
\midrule
\derived & computable from the primitives by theorem \\
\realunder & real, precise, predictive structure whose origin is not computable \\
\drawn & a coincidental draw of the abundance law \\
\bottomrule
\end{tabular}
\end{center}
The taxonomy classifies \emph{underivability}; \realunder\ and \drawn\ are \emph{not}
claimed as derivations.
\end{definition}

\section{The fourteen classified results}

\begin{center}
\begin{tabular}{llc}
\toprule
Object & Category & Confidence \\
\midrule
$U(1)$ & \derived & 0.95 \\
spatial 3 & \derived & 0.85 \\
$N\ge3$ & \derived & 0.90 \\
log-normal law & \derived & theorem \\
$SU(2)$ & \realunder\ (emergent) & 0.70 \\
$SU(3)$ structure & \realunder\ (emergent) & 0.10 \\
Koide $Q=2/3$ (reality) & \realunder\ (structured) & 0.90 \\
Koide $45^\circ$ (origin) & \realunder\ (structured) & 0.70 \\
Yukawas / couplings / $\Omega_{DM}$ & \drawn & — \\
$N\le3$ & \drawn & 0.70 \\
color count 3 & \drawn & — \\
internal $N=3$ & \derived$\cap$\drawn & 0.70 \\
neutrino-Koide & FALSIFIED & 0.90 \\
\bottomrule
\end{tabular}
\end{center}

Category conflicts: \textbf{0}. Composite objects: \textbf{2}. Overall consistency:
\textbf{0.95}.

%=============================================================================

\chapter{The No-Go Theorems}
\label{ch:nogo}

The no-go theorems are \emph{conditional} — relative to the no-new-primitives constraint and
the tested route space — not absolute logical impossibilities (except the computation
T-11).

\begin{center}
\begin{tabular}{lll}
\toprule
Theorem & Statement & Confidence \\
\midrule
\textbf{T-08} Koide no-go & no symmetry/attractor/topology/information-geometry selects $Q=2/3$ & 0.70 \\
\textbf{T-09} Gauge-count no-go (provisional) & no principle fixes $n$; graph-spectrum/lattice-mode untested & 0.10 \\
\textbf{T-10} $N\le3$ no-go & no stability/anomaly/representation bound gives $N\le3$ & 0.70 \\
\textbf{T-11} Neutrino-Koide falsification & $Q_{\max}=0.585/0.500 < 2/3$ (computation) & 0.90 \\
\textbf{T-12} Shared-cascade no-go & one vs three cascades untestable from one universe & 0.55 \\
\bottomrule
\end{tabular}
\end{center}

\begin{remark}[Internal-3 disposition]
Gauge count $n=3$ and multiplicity $N\le3$ are one node with two formally unlinked faces
($N$ vs $n$; assumption A-10). It is dispositioned \emph{unresolved-contingent}: the
multiplicity face closed (T-10, 0.70); the gauge-count face \emph{provisionally} closed
(T-09, 0.10). This is the theory's one open door.
\end{remark}

%=============================================================================

\chapter{Quantitative Predictions}
\label{ch:predictions}

\begin{center}
\begin{tabular}{lll}
\toprule
Prediction & Type & Status \\
\midrule
RAR $g_\dagger=cH_0/(2\pi)\approx1.05\times10^{-10}$ & zero-parameter & matches $a_0$ \\
$w(z)=-1+0.015(1+z)^{3/2}$ & cosmological & Euclid $>3\sigma$ by $\sim$2030 \\
$\Lambda(t)=\alpha/\sqrt{V(t)}$ & dark energy & Euclid / Roman \\
log-normal abundance law & distribution form & testable \\
$N\ge3$ (CP lower bound) & a priori & confirmed \\
neutrino-Koide $Q=2/3$ & falsifiable & \textbf{falsified} \\
\bottomrule
\end{tabular}
\end{center}

The theory is \emph{falsifiable}: it made a concrete prediction (neutrino-Koide) that was
falsified, and it carries live quantitative predictions that distinguish it from SM
$+\Lambda$CDM in the \emph{form} sector, even though the \emph{content} is DRAWN.

%=============================================================================

\chapter{Scope and Limitations}
\label{ch:scope}

\begin{enumerate}
\item The gauge-count face of the internal-3 node is only provisionally closed (T-09, 0.10).
\item Encyclopedia completeness $\approx81\%$; the CMB chapter is a documented partial
computational layer.
\item Content is contingent by design; the theory predicts the distribution \emph{shape},
not the \emph{values}.
\item The shared-cascade question is logically (not empirically) closed.
\item The unified action is a roadmap, not a result.
\item Confidence numbers are uncalibrated ordinal ranks (Phase-149/156 estimates).
\item The phase-gravity chain is ontological in the weak field.
\end{enumerate}

%=============================================================================
\part{Verification}
%=============================================================================

\chapter{The Executable Verification Program}
\label{ch:verification}

All claims below are backed by executable xUnit tests (all \textbf{PASSING}). The complete
inventory is in Appendix~\ref{app:tests}.

\section{Summary of verified results}

\begin{center}
\begin{tabular}{lll}
\toprule
Chain link & Test file & Result \\
\midrule
$L_Q\to-\nabla^2$ & GraphLaplacianContinuumTests & exact spectrum, 4$\times$/doubling \\
BDG$\to\square$ & BDGOperatorContinuumTests & $O(h^2)$, 4$\times$/halving \\
$L_Q$ vs BDG signature & QuantumGravityBridgeTests (3) & PSD vs indefinite \\
no native $\Delta_g$ & CurvedSpaceBridgeTests (3) & bridge absent \\
weighted Laplacian & WeightedLaplacianTests (4) & symmetric/PSD/zero-sum \\
$L_W\to\Delta_g$ & LaplaceBeltramiTests (3) & S¹ example \\
curved Schr\"odinger & CurvedSchrodingerTests (3) & unitary \\
Einstein recovery & EinsteinRecoveryTests (3) & PARTIAL \\
Einstein chain & EinsteinTensorTests (4) & standard 2D/3D \\
Einstein integration & EinsteinTensorIntegrationTests (4) & builder works \\
metric generation & MetricGenerationTests (4) & distance PRESENT \\
metric emergence & MetricEmergenceTests (4) & factor native \\
conformal structure & ConformalStructureTests (3) & class imported \\
metric origin & MetricOriginTests (3) & closure \\
\bottomrule
\end{tabular}
\end{center}

\section{The decisive findings}

\begin{enumerate}
\item The \emph{Schr\"odinger} side of the continuum limit is \textbf{controlled and
tested}: $L_Q\to-\nabla^2\to$ Schr\"odinger (and the curved $L_W$, and the
Laplace--Beltrami approximation) are verified.
\item The \emph{Einstein} side is \textbf{partially} verified: the standard $g\to
G_{\mu\nu}$ chain works, but TQM has no native metric (Malament import), and the BDG action
is imported.
\item The conformal-class ``gap'' is an \emph{imported theorem} (Malament), not a theory
gap.
\end{enumerate}

%=============================================================================

\chapter{The Hostile-Review Audit Trail}
\label{ch:audit-trail}

A hostile referee (Round~2) raised twelve FATAL issues. Each was re-evaluated against the
executable tests. The classifications are:

\begin{center}
\begin{tabular}{p{6.5cm}l}
\toprule
Round-2 FATAL issue & Classification \\
\midrule
Primitives undefined as mathematical objects & PARTIALLY RESOLVED \\
Schr\"odinger derivation circular/incomplete & RESOLVED \\
Gauge derivations = elementary topology & PARTIALLY RESOLVED \\
Complexity argument not a derivation & PARTIALLY RESOLVED \\
Structure/content split = immunization & PARTIALLY RESOLVED \\
Composite objects admitted & PARTIALLY RESOLVED \\
Internal-3 unresolved while ``complete'' & RESOLVED \\
T-09 at 0.10 cannot close & RESOLVED \\
``Structurally complete'' overstatement & RESOLVED \\
Gravity$\to$Einstein is ontological & RESOLVED \\
$\mathrm{Aut}(S^1)=U(1)$ as EM derivation & RESOLVED \\
No controlled continuum limit & PARTIALLY RESOLVED \\
\bottomrule
\end{tabular}
\end{center}

\textbf{Tally: RESOLVED = 6 · PARTIALLY RESOLVED = 6 · OPEN = 0.}

\section{What changed}

The framing overstatements (``closed'' vs.\ ``dispositioned'', ``no-go'' vs.\
``conditional no-go'', ``derivation'' vs.\ ``ontological reinterpretation'') were
corrected in the revised paper. The \emph{substantive} objection — ``no controlled
continuum limit'' — was met on the Schr\"odinger side (now tested: $L_Q\to-\nabla^2\to$
Schr\"odinger, and curved $L_W$, Laplace--Beltrami, unitary), and \emph{partially} on the
Einstein side (the standard chain works, but the metric and BDG action are imported).

\section{The decisive residual}

The single genuinely-open item is the \textbf{native metric $\to$ operator coupling} (the
``G4'' gap): TQM imports $g_{\mu\nu}$ (Malament) rather than generating it from Q-events.
The metric \emph{origin} is closed (Chapter~\ref{ch:metric-emergence}); the metric
\emph{dynamics} remain imported.

\textbf{Verdict:} READY\_FOR\_WHITEPAPER — NOT\_READY\_FOR\_JOURNAL (see
Chapter~\ref{ch:conclusion}).

%=============================================================================

\chapter{Conclusion}
\label{ch:conclusion}

THE Q-MODEL is \emph{structurally complete} in the precise sense that every in-scope
question is \textbf{dispositioned} — derived, underivable, underdetermined, contingent, or
accepted as a computational layer.

The single residual of genuine scientific tension is the \textbf{internal-3 node} — why the
internal multiplicity/count saturates at 3 — dispositioned unresolved-contingent with a
provisional gauge-count no-go (T-09, 0.10). This is the theory's one open door.

\textbf{Structure is derivable. Content is realized.}

\textbf{Publication status:} READY\_FOR\_WHITEPAPER — NOT\_READY\_FOR\_JOURNAL. The
Schr\"odinger dynamical core is executable and tested; the Einstein recovery remains
logical, not mathematical (imported metric + imported BDG action, $G$ dimensional, no
unique discriminating prediction).

%=============================================================================
\appendix
%=============================================================================

\chapter{Complete Test Inventory}
\label{app:tests}

\begin{longtable}{lllc}
\toprule
Test file & Tests & Verified quantity & Status \\
\midrule
\endhead
GraphLaplacianContinuumTests & 1 & $\lambda_k=(1/\Delta x^2)[2-2\cos(\pi k/(N{+}1))]\to(\pi k)^2$ & PASS \\
BDGOperatorContinuumTests & 1 & BDG$\to\square$, $O(h^2)$ & PASS \\
QuantumGravityBridgeTests & 3 & $L_Q$ PSD; BDG indefinite; incompatible & PASS \\
CurvedSpaceBridgeTests & 3 & no $\Delta_g$; $\Delta_g\to\nabla^2$; absent & PASS \\
MetricOperatorTests & 4 & candidate constructions & PASS \\
WeightedLaplacianTests & 4 & symmetric; zero-sum; PSD; unweighted limit & PASS \\
LaplaceBeltramiTests & 3 & flat limit; weights; S¹ ($2k{-}1$ degeneracy) & PASS \\
CurvedSchrodingerTests & 3 & unitary; flat limit; curved operator & PASS \\
EinsteinRecoveryTests & 3 & PARTIAL (no full $G_{\mu\nu}$) & PASS \\
EinsteinTensorTests & 4 & flat/sph $\Gamma$, $K$, $R$, $G$ & PASS \\
EinsteinTensorIntegrationTests & 4 & builder; 3-sphere $G=-g$ & PASS \\
MetricGenerationTests & 4 & distance PRESENT; candidate PARTIAL & PASS \\
MetricEmergenceTests & 4 & factor $f=\rho^{2/d}$; $R=-0.6145$ & PASS \\
ConformalStructureTests & 3 & axioms; null invariant; Malament & PASS \\
MetricOriginTests & 3 & class unique; factor free; imported & PASS \\
\bottomrule
\end{longtable}

\chapter{Confidence and Classification Tables}
\label{app:confidence}

The confidence numbers are uncalibrated ordinal ranks (Phase-149/156 estimates), not
posterior probabilities. They rank, they do not measure.

\section{Confidence table}

\begin{center}
\begin{tabular}{llc}
\toprule
Result & Confidence & Note \\
\midrule
program consistency & 0.95 & zero category conflicts \\
$U(1)$ & 0.95 & theorem \\
$N\ge3$ & 0.90 & CP lower bound \\
neutrino-Koide & 0.90 & falsification (computation) \\
spatial 3 & 0.85 & window intersection \\
$SU(2)$ & 0.70 & emergent \\
Koide $45^\circ$ & 0.70 & structured \\
$N\le3$ & 0.70 & no-go \\
$SU(3)$ structure & 0.10 & provisional no-go \\
\bottomrule
\end{tabular}
\end{center}

\section{Final classification}

\begin{center}
\begin{tabular}{ll}
\toprule
Question & Classification \\
\midrule
Is the metric origin closed? & YES (imported Malament, proven) \\
Is the Schr\"odinger continuum limit controlled? & YES (tested) \\
Is the Einstein recovery a derivation? & NO (logical, imported) \\
Is TQM journal-ready? & NO \\
Is TQM white-paper-ready? & YES \\
\bottomrule
\end{tabular}
\end{center}

%=============================================================================

\chapter{Worked Derivation: The Einstein Chain on the 2-Sphere}
\label{app:worked}

This appendix exhibits the standard differential-geometry chain explicitly, so the reader
can follow every step without opening the repository code.

\section{Setup}

The unit 2-sphere has metric (coordinates $x^1=\theta$, $x^2=\phi$)
\begin{equation}
g = \mathrm{diag}\!\left(g_{\theta\theta},\ g_{\phi\phi}\right)
  = \mathrm{diag}\!\left(1,\ \sin^2\theta\right),
\qquad g^{\theta\theta}=1,\quad g^{\phi\phi}=\frac{1}{\sin^2\theta}.
\end{equation}

\section{Christoffel symbols}

\begin{equation}
\Gamma^\lambda_{\mu\nu} = \tfrac12 g^{\lambda\sigma}\left(
\partial_\mu g_{\sigma\nu} + \partial_\nu g_{\sigma\mu} - \partial_\sigma g_{\mu\nu}\right).
\end{equation}

The only non-vanishing partial derivative is
$\partial_\theta g_{\phi\phi} = 2\sin\theta\cos\theta$. Therefore
\begin{align}
\Gamma^\theta_{\phi\phi} &= -\tfrac12\, g^{\theta\theta}\, \partial_\theta g_{\phi\phi}
= -\sin\theta\cos\theta, \\
\Gamma^\phi_{\theta\phi} = \Gamma^\phi_{\phi\theta}
&= \tfrac12\, g^{\phi\phi}\, \partial_\theta g_{\phi\phi}
= \frac{\cos\theta}{\sin\theta} = \cot\theta.
\end{align}
At $\theta=\pi/4$: $\Gamma^\theta_{\phi\phi}=-1/2$, $\Gamma^\phi_{\theta\phi}=1$.

\section{Riemann curvature}

\begin{equation}
R^\rho_{\ \sigma\mu\nu} = \partial_\mu\Gamma^\rho_{\nu\sigma}
- \partial_\nu\Gamma^\rho_{\mu\sigma}
+ \Gamma^\rho_{\mu\lambda}\Gamma^\lambda_{\nu\sigma}
- \Gamma^\rho_{\nu\lambda}\Gamma^\lambda_{\mu\sigma}.
\end{equation}

The single independent component is
\begin{align}
R^\theta_{\ \phi\theta\phi}
&= \partial_\theta\Gamma^\theta_{\phi\phi} - \Gamma^\theta_{\phi\phi}\Gamma^\phi_{\theta\phi} \\
&= -\left(\cos^2\theta-\sin^2\theta\right) + \sin\theta\cos\theta\cdot\cot\theta \\
&= \sin^2\theta.
\end{align}
At $\theta=\pi/4$: $R^\theta_{\ \phi\theta\phi}=1/2$.

\section{Gauss curvature, Ricci, and Einstein}

Gauss curvature:
\begin{equation}
K = \frac{R^\theta_{\ \phi\theta\phi}}{g_{\phi\phi}}
= \frac{\sin^2\theta}{\sin^2\theta} = 1.
\end{equation}

Ricci tensor (in 2D, $R_{\mu\nu}=K g_{\mu\nu}$):
\begin{equation}
R_{\theta\theta}=K g_{\theta\theta}=1,\qquad
R_{\phi\phi}=K g_{\phi\phi}=\sin^2\theta=\tfrac12 \text{ at }\theta=\pi/4.
\end{equation}

Ricci scalar:
\begin{equation}
R = g^{\mu\nu}R_{\mu\nu} = 2K = 2.
\end{equation}

Einstein tensor (vanishes identically in 2D):
\begin{equation}
G_{\mu\nu} = R_{\mu\nu} - \tfrac12 R g_{\mu\nu} = 0.
\end{equation}

All of these are reproduced numerically by \texttt{EinsteinTensorBuilder} to within
$10^{-3}$. A non-trivial $G_{\mu\nu}$ requires dimension $\ge3$; the unit 3-sphere gives
$G_{\mu\nu}=-g_{\mu\nu}$ (Appendix~\ref{app:tests}).

%=============================================================================

\chapter{Research Programs}
\label{app:programs}

The derivation program is organized into research programs, each closed by executable
experiments:

\begin{center}
\begin{tabular}{llc}
\toprule
Program & Scope & Experiments \\
\midrule
DATA & Cosmology \& RAR & 10 (DATA-001$\to$010) \\
QM & Quantum foundations & 5 (QM-001$\to$005) \\
QG & Quantum gravity & 31 (QG-001$\to$031) \\
XC & Continuum-limit bridge & (audits + tests) \\
\bottomrule
\end{tabular}
\end{center}

Key results of the full program:
\begin{itemize}
\item $G$ is not fundamental: $G=\ell^2c^3/\hbar$ (gravity emerges from the triple).
\item $(\ell,\tau,\hbar)$ is the irreducible triple — one process, three aspects.
\item Oscillation is a logical inevitability (cannot be removed without destroying $Q$).
\item Gravity is a phase-gradient phenomenon; attraction dominates, repulsion is dark
energy.
\item Gravity manipulation is not possible (QG-023$\to$031).
\item RAR derivation $g_\dagger=cH_0/(2\pi)$ with zero free parameters.
\item Parameter compression SM+GR ($\sim$26) $\to$ TQM ($2{+}3$), $\sim$5--8$\times$
reduction.
\item Novel predictions: $w(z)=-1+0.015(1+z)^{3/2}$; evolving $g_\dagger(z)$; Euclid
sensitivity.
\end{itemize}

\chapter{Glossary}
\label{app:glossary}

\begin{center}
\begin{tabular}{lp{10cm}}
\toprule
Term & Meaning \\
\midrule
$Q$ & principle of individuation (ontology) \\
Random Actualization & ontological chance (assumption A-03) \\
$(\ell,\tau,\hbar)$ & spacetime scale, clock, action (units) \\
$M^2$ & nonlinearity parameter (single continuous) \\
$L_Q$ & graph Laplacian $D-A$ \\
$L_W$ & weighted Laplacian $D_K-K$ \\
BDG & Benincasa--Dowker--Glaser d'Alembertian \\
$\Delta_g$ & Laplace--Beltrami operator \\
$G_{\mu\nu}$ & Einstein tensor $R_{\mu\nu}-\tfrac12 R g_{\mu\nu}$ \\
\derived & computable from primitives by theorem \\
\realunder & real structure, underivable origin \\
\drawn & coincidental draw of the abundance law \\
T-08\dots T-12 & conditional no-go theorems \\
A-03, A-06, A-07, A-10 & assumptions (random actualization; $Z_2$ lift; gauge dynamics; $N$/$n$ split) \\
\bottomrule
\end{tabular}
\end{center}

%=============================================================================

\chapter{References}
\label{app:references}

\begin{thebibliography}{99}

\bibitem{malament1977}
D.~B.~Malament, ``The class of continuous timelike curves determines the topology of
spacetime,'' \emph{J. Math. Phys.} \textbf{18}, 1399 (1977).

\bibitem{hkm1976}
S.~W.~Hawking, A.~R.~King, P.~J.~McCarthy, ``A new topology for curved space--time which
incorporates the causal, differential, and conformal structures,'' \emph{J. Math. Phys.}
\textbf{17}, 174 (1976).

\bibitem{myrheim1978}
J.~Myrheim, ``Statistical geometry,'' CERN-TH-2538 (1978).

\bibitem{bdg2010}
D.~M.~T.~Benincasa, F.~Dowker, ``Scalar curvature of a causal set,'' \emph{Phys. Rev. Lett.}
\textbf{104}, 181301 (2010); L.~Glaser, \emph{Class. Quant. Grav.} \textbf{31}, 095007
(2014).

\bibitem{gleason1957}
A.~M.~Gleason, ``Measures on the closed subspaces of a Hilbert space,''
\emph{J. Math. Mech.} \textbf{6}, 885 (1957).

\bibitem{sorkin}
R.~D.~Sorkin, ``Causal sets: discrete gravity,'' \emph{Proceedings of the Valdivia Summer
School} (2003).

\bibitem{belkin2003}
M.~Belkin, P.~Niyogi, ``Laplacian eigenmaps for dimensionality reduction and data
representation,'' \emph{Neural Computation} \textbf{15}, 1373 (2003).

\bibitem{koide1981}
Y.~Koide, ``Fermion mass relation and the generation structure,'' \emph{Phys. Rev. D}
\textbf{28}, 252 (1983) [letter 1981].

\bibitem{bertrand}
J.~Bertrand, ``Th\'eor\`eme relatif au mouvement d'un point attir\'e vers un centre fixe,''
\emph{C. R. Acad. Sci.} \textbf{77}, 849 (1873).

\end{thebibliography}

\end{document}
